\begin{table}[t]
\centering
\caption{Runtime settings (median ms). Turning garbage collection off (Go) or making it rare (\cs{}) has no measurable effect; disabling tiered compilation makes short \cs{} runs 2--4$\times$ faster.}
\label{tab:runtime}
\small
\begin{tabular}{lrrrrrr}
\toprule
 & \multicolumn{2}{c}{Eppstein $200^2$, $k=10^4$} & \multicolumn{2}{c}{Eppstein $1000^2$, $k=10^5$} & \multicolumn{2}{c}{Yen $200^2$, $k=10^3$} \\ \cmidrule(lr){2-3}\cmidrule(lr){4-5}\cmidrule(lr){6-7}
 & tree & next & tree & next & tree & next \\
\midrule
\cs{} default & 27.7 & 76.2 & \num{255} & \num{3441} & 38.1 & \num{5485} \\
\cs{} gen0 budget 256\,MB & 28.5 & 78.9 & \num{257} & \num{3489} & 38.4 & \num{5428} \\
\cs{} without tiered JIT & 6.4 & 42.9 & \num{248} & \num{3332} & 8.6 & \num{5331} \\
Go default & 9.0 & 43.9 & \num{324} & \num{3468} & 10.9 & \num{7581} \\
Go \texttt{GOGC=off} & 8.0 & 42.7 & \num{318} & \num{3451} & 12.4 & \num{7646} \\
Rust & 6.6 & 34.8 & \num{259} & \num{3092} & 9.4 & \num{5869} \\
\bottomrule
\end{tabular}
\end{table}


\Cref{tab:runtime} separates the effects of the runtimes. \emph{Garbage collection is not a factor here.} Turning it
off completely in Go (\texttt{GOGC=off}) or making it rare in \cs{} (a 256\,MB generation-0 budget) changes no
measured time beyond noise. The reason is the data layout, not the collectors: all per-vertex data lives in arrays
that are reused across queries, and the only allocations during an enumeration are small, short-lived candidate and
heap objects, which generational collectors handle cheaply.

\emph{JIT warm-up is a factor.} With default settings, .NET first compiles methods quickly without optimization
(tier~0) and optimizes them only after many calls or long loops. A benchmark that runs a query twice therefore
measures partly unoptimized code. With tiered compilation disabled, the \cs{} tree on the $200^2$ grid takes 6.4\,ms
instead of 27.7\,ms (Go 9.0\,ms, Rust 6.6\,ms), and Yen's tree 8.6\,ms instead of 38.1\,ms. For long-running
services this does not matter; for command-line tools and short benchmarks it dominates.

\emph{What remains is compilers and libraries.} Rust is fastest in most scenarios. Its largest advantages are in
building the graph (253\,ms by index versus 440 and 530\,ms) and in path enumeration (10--30\%), where arenas with
32-bit links replace small objects. In the hot loops of Dijkstra and A*, all three languages do the same array
accesses with bounds checks, and their times are within about 25\% when \cs{} runs optimized code. Library choices
also matter at this level: the 4-ary heap of .NET versus our binary heaps, and Go maps keyed by strings versus a
\cs{} hash set or a Rust hash set.

\emph{Memory.} The live heap of the $10^6$-vertex graph with one query is similar: 117\,MB (\cs), 94\,MB (Go) and
112\,MB (Rust). Go and \cs{} share the caller's label strings, which are therefore not counted. The live heap
understates what a garbage-collected process needs, because collectors run with headroom~\cite{hertz2005}; we did
not measure peak resident memory.

\emph{Engineering effort.} The Rust port was the third implementation and was written by the agent from the Go
version. Two Rust-specific properties stood out. The borrow checker turned the ``graph changed during iteration''
error of \cs{} and Go into a compile-time error, at the price of a slightly less convenient iterator. And the most
natural Rust structure for persistent heaps, \texttt{Rc} pointers, was replaced by index-based arenas, which are
both faster and simpler to reason about than reference counting.
